\documentclass[11pt,a4paper]{article}

% --- Packages ---
% Target: Royal Society Open Science (open access, interdisciplinary,
% published Gero et al. 2016; explicitly welcomes cross-disciplinary work)
% Alternative targets: Bioacoustics, PLOS ONE, JASA, NeurIPS workshop
\usepackage[utf8]{inputenc}
\usepackage[T1]{fontenc}
\usepackage{amsmath,amssymb,amsfonts}
\usepackage{graphicx}
\usepackage{booktabs}
\usepackage{hyperref}
\usepackage[margin=2.5cm]{geometry}
\usepackage{natbib}
\usepackage{xcolor}
\usepackage{caption}
\usepackage{subcaption}
\usepackage{authblk}
\usepackage{lineno}
\linenumbers

% --- Title ---
\title{Geometric Structure in Sperm Whale Communication:\\
Hyperbolic Embeddings, Topological Analysis, and Adversarial Robustness}

\author[1]{Andrew H.\ Bond}
\affil[1]{Department of Computer Engineering, San Jos\'e State University, San Jos\'e, CA, USA}

\date{}

\begin{document}
\maketitle

% ===================================================================
\begin{abstract}
Recent work has revealed that sperm whale (\textit{Physeter macrocephalus})
codas possess a combinatorial phonetic structure comparable in complexity
to human phonological systems. We apply differential geometry, algebraic
topology, and adversarial robustness testing to 8{,}719 annotated codas
from the Dominica Sperm Whale Project. We show that: (1)~the hierarchical
coda type taxonomy embeds with lower distortion in the Poincar\'e ball
than in Euclidean space, consistent with tree-like combinatorial structure;
(2)~persistent homology on inter-click interval point clouds reveals
topologically distinct attractors per rhythm class;
(3)~adversarial acoustic perturbations map asymmetric phonetic decision
boundaries (mean asymmetry $= 0.87$), analogous to directional vowel
chain shifts in human language; (4)~the coda system satisfies
Menzerath's law ($r = -0.269$, $p < 10^{-144}$), exhibits higher-order
Markov sequential structure, and supports active turn-taking with
cross-whale response latency approximately half that of same-whale
continuation; and (5)~individual whales produce the same coda type with
statistically distinguishable inter-click interval patterns
(Kolmogorov--Smirnov $p < 0.001$ for all tested pairs), confirming
individual identity encoding. We introduce the Decoder Robustness Index
(DRI), the first adversarial robustness benchmark for cetacean
communication decoders, and demonstrate that adversarial perturbation
serves as a phonetic boundary discovery tool. All methods are released
as the open-source \texttt{eris-ketos} toolkit.
\end{abstract}

\noindent\textbf{Keywords:} sperm whale communication, hyperbolic
embeddings, topological data analysis, adversarial robustness,
computational bioacoustics, linguistic universals

% ===================================================================
\section{Introduction}

Sperm whales produce stereotyped sequences of broadband clicks known as
codas, which serve as communicative signals mediating social identity and
group cohesion~\citep{Weilgart1993,Rendell2003}. Individual codas
consist of 3--12 clicks with characteristic inter-click intervals (ICIs)
that encode coda type. Social units within vocal clans share coda
repertoires but exhibit unit-level and individual-level variation in
production~\citep{Gero2016,Hersh2022}.

A landmark study by \citet{Sharma2024} demonstrated that sperm whale
codas possess a combinatorial phonetic structure decomposable into four
independently controlled features: rhythm (18 types), tempo (5 types),
rubato (3 types), and ornamentation (2 types), yielding a combinatorial
space of up to 540 configurations. Subsequent work identified
vowel-like spectral patterns within individual clicks, adding a fifth
phonetic dimension~\citep{Begus2025}. These findings elevate the
complexity of cetacean communication systems to a level that invites
methods from differential geometry, algebraic topology, and adversarial
machine learning---tools that have proven effective for analysing
hierarchical, manifold-valued, and combinatorial structure in other
domains but have not previously been applied to animal communication.

Three specific properties of the coda system motivate geometric
methods. First, the combinatorial taxonomy of coda types is inherently
tree-like (rhythm class $\to$ click count $\to$ variant), and tree-like
structures embed with $O(\log n)$ distortion in hyperbolic space versus
$O(n)$ in Euclidean space~\citep{Sarkar2011}. The Poincar\'e ball
model~\citep{Nickel2017} is therefore a natural representation for the
coda type hierarchy. Second, inter-click interval patterns within each
coda type form point clouds in $\mathbb{R}^k$ whose topology---connected
components, loops, voids---may encode structural properties of the
communication system that summary statistics miss. Persistent
homology~\citep{Carlsson2009,Bauer2021} provides a principled framework
for extracting such topological features. Third, the robustness of coda
classification to acoustic perturbation is an open question with
implications for both decoder design and communication theory: if small
perturbations flip classification, the system has low redundancy and
high information density.

In this work, we present five contributions:

\begin{enumerate}
    \item \textbf{Poincar\'e ball embedding} of the coda type hierarchy,
    with quantitative distortion comparison against Euclidean baselines.

    \item \textbf{Persistent homology} on ICI point clouds, revealing
    topologically distinct attractors per rhythm class.

    \item \textbf{Decoder Robustness Index (DRI)}, the first adversarial
    robustness benchmark for cetacean communication decoders, adapted
    from the Bond Index framework for ethical AI evaluation.

    \item \textbf{Adversarial perturbation as a discovery tool},
    revealing asymmetric phonetic boundaries and empirically estimating
    channel capacity at approximately 3.0 bits per coda.

    \item \textbf{Linguistic universals}, confirming Menzerath's law,
    higher-order Markov structure, individual accents, and active
    turn-taking in coda exchanges.
\end{enumerate}

All methods are implemented in the open-source \texttt{eris-ketos}
Python package (\texttt{pip install eris-ketos}) and demonstrated on the
full Dominica Sperm Whale Project (DSWP)
dataset~\citep{Sharma2024}.

% ===================================================================
\section{Results}

\subsection{Poincar\'e ball embedding captures hierarchical structure}

We parsed the 28 coda types in the DSWP dataset into a three-level
taxonomy (rhythm class $\to$ click count $\to$ variant) and constructed
a taxonomic distance matrix. Coda types were embedded into a
$d$-dimensional Poincar\'e ball ($d = \min(16, n_\mathrm{types})$) via
gradient descent on a distortion loss, and compared against Euclidean
embedding via PCA on the same distance matrix (Fig.~\ref{fig:poincare}).

The Poincar\'e embedding achieved higher Spearman rank correlation
($\rho$) between embedded pairwise distances and true taxonomic
distances than the Euclidean baseline, consistent with the theoretical
advantage of hyperbolic geometry for tree-like structures. A Hyperbolic
Multinomial Logistic Regression classifier trained on individual coda
feature vectors (9 normalised ICI features, log-duration, click count)
achieved competitive classification accuracy against Euclidean baselines
(logistic regression, $k$-nearest neighbours, random forest), with class
prototypes initialised from the taxonomic Poincar\'e embeddings.

\subsection{Topological data analysis reveals distinct rhythm attractors}

For each coda type with $\geq 40$ samples, we constructed an ICI point
cloud in $\mathbb{R}^9$ and computed Vietoris--Rips persistent homology
up to dimension 1 using Ripser~\citep{Bauer2021,Tralie2018}. The
resulting persistence diagrams (Fig.~\ref{fig:tda}) revealed
systematic differences across rhythm classes:

\begin{itemize}
    \item \textbf{Regular codas} (xR) exhibited tight clusters with low
    $H_1$ persistence, reflecting their temporally uniform ICI patterns.

    \item \textbf{Irregular codas} (xi) produced higher $H_1$
    persistence, indicating topological loops in the ICI space consistent
    with greater timing variability.

    \item \textbf{Compound codas} (x+y) showed elevated $H_0$
    persistence, reflecting multi-component structure.
\end{itemize}

These topological signatures are invariant to monotone transformations of
the distance metric and thus capture structural properties that
Euclidean summary statistics cannot.

\subsection{Adversarial perturbation maps phonetic boundaries}

We applied nine parametric acoustic transforms (additive noise, pink
noise, Doppler shift, multipath echo, amplitude scaling, time stretch,
spectral mask, click dropout, time shift) at 11 graduated intensity
levels (0.0 to 1.0) to synthesised coda signals and measured boundary
crossing rates in ICI feature space (Fig.~\ref{fig:fuzz}).

\paragraph{Transform sensitivity spectrum.}
Click dropout caused the highest boundary crossing rate (50\% at full
intensity), confirming that individual clicks are the primary
information-carrying elements. Time shift, despite being theoretically
invariant for a correct decoder, caused 40\% crossings, exposing decoder
vulnerability. Multipath echo and spectral mask caused $<2\%$ crossings,
suggesting these acoustic dimensions carry minimal semantic content in
the ICI-based representation.

\paragraph{Phonetic boundary asymmetry.}
The phonetic neighbourhood map revealed a mean boundary asymmetry of
0.87 (on a 0--1 scale where 0 indicates symmetric and 1 indicates
unidirectional transitions). For example, 5R3 $\to$ 4R2 exhibited 223
crossings while 4R2 $\to$ 5R3 exhibited 0. This asymmetry is
reminiscent of directional vowel chain shifts in human
language~phonology, where perturbation-induced transitions favour
specific directions.

\paragraph{Per-type robustness.}
Irregular codas exhibited the highest mean activation threshold (0.944),
while compound codas were least robust (0.800). This ordering suggests
that temporal irregularity occupies a wider phonetic territory in
acoustic space.

\paragraph{Channel capacity estimate.}
All eight tested coda types remained well-separated under perturbation
(activation threshold $> 0.5$), yielding an empirical channel capacity
estimate of $\log_2(8) = 3.0$ bits per coda. This exceeds the unigram
entropy of the rhythm distribution alone ($H(R) = 2.1$ bits), consistent
with multi-feature information encoding.

\subsection{Linguistic universals in coda communication}

Analysis of the full DSWP dataset and the associated exchange-level
dialogues data revealed multiple statistical universals of natural
language (Fig.~\ref{fig:linguistic}):

\paragraph{Menzerath's law.}
Longer codas (more clicks) had shorter mean ICIs (Spearman $r = -0.269$,
$p < 10^{-144}$, $n = 8{,}719$), confirming the linguistic universal
that longer constructs are composed of shorter constituents. This result
is consistent with the cross-species meta-analysis
of~\citet{Youngblood2025}.

\paragraph{Zipf's law.}
The rank-frequency distribution of coda types followed a power law with
exponent $\alpha = 1.44$, steeper than the $\alpha \approx 1.0$ typical
of human language, indicating a more skewed distribution with a few
dominant types (1+1+3, 5R1, 5R3).

\paragraph{Sequential structure.}
Coda exchanges exhibited first-order Markov dependencies: the mutual
information between consecutive rhythm types was $I(R_t; R_{t+1}) =
0.380$ bits, representing 18.0\% of the unigram entropy. Trigram
prediction accuracy (71.2\%) exceeded bigram accuracy (66.1\%),
demonstrating higher-order sequential grammar.

\paragraph{Turn-taking.}
Cross-whale response latency ($\bar{x} = 2.25$~s, $n = 2{,}225$) was
approximately half that of same-whale continuation latency ($\bar{x} =
4.68$~s, $n = 1{,}379$), with highly significant distributional
difference (Kolmogorov--Smirnov $D = 0.560$, $p < 10^{-247}$). This
indicates active turn-taking rather than independent vocalisation.

\paragraph{Individual accents.}
For the most common coda type (1+1+3), pairwise Kolmogorov--Smirnov
tests between ICI distributions of four identified whales revealed
statistically significant differences for all six pairs ($p < 0.001$),
confirming that individuals encode identity information within
supposedly standardised coda types~\citep{Gero2016,Hersh2022}.

% ===================================================================
\section{Discussion}

The results presented here demonstrate that geometric and topological
methods reveal structure in cetacean communication that standard
Euclidean analyses fail to capture. The Poincar\'e ball embedding
naturally represents the hierarchical taxonomy of coda types, persistent
homology distinguishes rhythm classes by their topological signatures,
and adversarial perturbation maps phonetic decision boundaries with a
precision that traditional classification accuracy cannot achieve.

The discovery of highly asymmetric phonetic boundaries (mean asymmetry
$= 0.87$) is particularly noteworthy. In human phonology, directional
chain shifts---where perturbation-induced transitions favour specific
directions---are a well-studied phenomenon that reflects the geometry of
the phonetic space. The observation of analogous asymmetry in sperm
whale codas suggests that the coda combinatorial space may possess a
non-trivial geometric structure with preferred transition directions,
potentially reflecting acoustic or perceptual constraints on coda
production and reception.

The confirmation of multiple linguistic universals (Menzerath's law,
Zipf's law, higher-order Markov structure, active turn-taking) adds to
the growing body of evidence~\citep{Youngblood2025} that cetacean
communication systems share deep statistical regularities with human
language. The finding that individual whales encode identity within
ostensibly standardised coda types extends the
identity-encoding hierarchy documented by \citet{Gero2016}
and~\citet{Hersh2022}.

The Decoder Robustness Index introduced here provides the first
systematic framework for evaluating the reliability of cetacean
communication decoders under realistic acoustic perturbations. The
finding that click dropout is the most disruptive transform (50\%
crossing rate) while spectral masking is nearly harmless ($<1\%$)
has direct implications for decoder design: robustness to missing clicks
should be prioritised over spectral invariance.

\paragraph{Limitations.}
The present analysis relies on ICI-based features from the annotated
DSWP dataset. Extension to raw acoustic features---including the
vowel-like spectral patterns reported by \citet{Begus2025}---would
require waveform-level analysis with real recordings. The DRI
measurements use a simple nearest-centroid decoder; evaluation of more
sophisticated models such as WhAM~\citep{Paradise2025} would provide a
more ecologically meaningful robustness benchmark. The channel capacity
estimate of 3.0 bits per coda is based on the number of
well-separated types under perturbation and should be regarded as a
lower bound.

\paragraph{Future directions.}
Three avenues are most promising. First, application of the geometric
framework to other cetacean species---particularly orcas (DCLDE dataset,
40+ stereotyped call types per community) and humpback whales
(hierarchical song grammar)---to test whether communication systems that
evolved independently converge on similar geometric structures. Second,
integration of exchange-level sequential structure with geometric
embeddings to model the temporal dynamics of coda conversations as
trajectories on manifolds. Third, collaboration with field researchers
(Project CETI, DSWP) to validate geometric predictions against
behavioural context and to design playback experiments testing whether
synthetic codas generated from geometric models elicit appropriate
responses.

% ===================================================================
\section{Methods}

\subsection{Data}
We used the annotated coda dataset from the Dominica Sperm Whale
Project~\citep{Sharma2024}, comprising 8{,}719 codas with inter-click
interval measurements, coda type labels, social unit identifiers, and
individual whale identifiers. The exchange-level dialogues dataset
(3{,}840 entries) provided sequential ordering and rhythm type labels for
Markov and turn-taking analyses. Both datasets are publicly available
under CC BY 4.0 licence.

\subsection{Poincar\'e ball embedding}
Coda types were parsed into a three-level taxonomy (rhythm class $\to$
click count $\to$ variant). A taxonomic distance matrix was computed
as the number of levels at which two types first differ (0--4). Types
were embedded into a Poincar\'e ball of curvature $c = 1.0$ and
dimension $d = \min(16, n_\mathrm{types})$ via gradient descent
minimising the mean squared error between embedded hyperbolic distances
and true taxonomic distances. Euclidean baselines used PCA on the same
distance matrix.

\subsection{Classification}
A Hyperbolic Multinomial Logistic Regression (HyperbolicMLR) classifier
was trained on 11-dimensional feature vectors (9 duration-normalised ICI
features, log-duration, normalised click count) with class prototypes
initialised from the taxonomic embeddings. Features were standardised
and mapped to the Poincar\'e ball via the exponential map at the origin,
with scaling calibrated to the 95th percentile of input norms. Training
used Adam optimisation with cosine learning rate annealing (500 epochs,
initial LR $= 0.02$, minimum LR $= 10^{-4}$). Euclidean baselines
(logistic regression, $k$-NN with $k = 5$, random forest with 100
trees) were trained on the same features.

\subsection{Topological data analysis}
For each coda type with $\geq 40$ samples, an ICI point cloud was
constructed in $\mathbb{R}^9$. Point clouds exceeding 300 points were
subsampled with fixed random seed. Vietoris--Rips persistent homology
was computed to dimension 1 using Ripser~\citep{Bauer2021}. Features
extracted from persistence diagrams included $H_0$ and $H_1$ counts,
maximum and total lifetime, mean birth and death times, and entropy of
the persistence distribution.

\subsection{Decoder Robustness Index}
Nine parametric acoustic transforms were applied at 11 intensity levels
(0.0 to 1.0) to synthesised coda signals. Coda signals were synthesised
from ICI data as Gaussian-enveloped 5~kHz click trains. The DRI was
computed as $\text{DRI} = 0.5 \cdot \bar{\omega} + 0.3 \cdot
\omega_{75} + 0.2 \cdot \omega_{95}$, where $\omega$ is the graduated
semantic distance between baseline and perturbed classifications.
Adversarial thresholds were found via binary search for the minimal
intensity causing a classification flip.

\subsection{Statistical analyses}
Zipf's law was assessed by fitting a power law $f(r) = C \cdot
r^{-\alpha}$ to the rank-frequency distribution. Menzerath's law was
tested via Spearman rank correlation between click count and mean ICI.
Sequential structure was quantified via unigram entropy, bigram
conditional entropy, and mutual information. Individual accents were
tested via pairwise Kolmogorov--Smirnov tests on ICI distributions.
Turn-taking was assessed by comparing inter-coda intervals for same-whale
versus cross-whale transitions within recording sessions.

\subsection{Software}
All analyses were implemented in the \texttt{eris-ketos} Python package
(v0.1.0), available at \url{https://pypi.org/project/eris-ketos/} and
\url{https://github.com/ahb-sjsu/eris-ketos}. The package depends on
NumPy, PyTorch, scikit-learn, Ripser, and librosa. A self-contained
Jupyter notebook reproducing all results is included in the repository.

% ===================================================================
\section*{Data availability}
The DSWP coda annotation dataset is available at
\url{https://github.com/pratyushasharma/sw-combinatoriality} under
CC BY 4.0 licence. Raw audio recordings are available on HuggingFace
at \url{https://huggingface.co/datasets/orrp/DSWP}.

\section*{Code availability}
The \texttt{eris-ketos} package is available at
\url{https://github.com/ahb-sjsu/eris-ketos} under MIT licence.

\section*{Acknowledgements}
We thank Shane Gero and the Dominica Sperm Whale Project for making the
coda annotation data publicly available, and Pratyusha Sharma and
colleagues for releasing the combinatorial analysis code and data. The
geometric analysis toolkit builds on the ErisML framework for
structured ethical evaluation.

% ===================================================================
\bibliographystyle{plainnat}
\bibliography{references}

\end{document}
